% ============================================================
% §5  Discussion
% ============================================================
\section{Discussion}
\label{sec:discussion}

\subsection{Exploitation Is the Critical Ingredient}
\label{sec:exploitation}

The most consequential result is not that EI and UCB outperform Random—
that is expected—but that Uncertainty is statistically indistinguishable
from Random ($p = 0.974$, average ranks $3.38$ vs.\ $3.48$).
Uncertainty is not a blind strategy: it fits the same GP as EI and UCB
at every iteration and selects the point of maximum posterior standard
deviation.
Yet its rank matches that of blind Sobol sampling across 60
(problem, seed) blocks.

This isolates where the value of BO lies.
Uncertainty and Random share one characteristic: neither uses the GP's
posterior mean.
EI and UCB share the complementary characteristic: both combine the
posterior mean with the posterior variance.
The equivalence classes therefore provide a clean experimental
decomposition—the GP's uncertainty estimate alone is insufficient;
exploitation of the mean is the necessary component.

Two independent lines of theory predict this outcome.
De~Ath et al.~\citeyear{deAth2019greed} recast acquisition function
selection as a bi-objective problem: they show that EI and UCB both
select points on the exploration–exploitation Pareto front, whereas a
strategy that uses only $\sigma(x)$ (the pure-exploration extreme) sits
off this front in the exploitation dimension.
Shahriari et al.~\citeyear{shahriari2016taking} observe that acquisition
functions matter, but caution against over-indexing on any single one;
their framing implies that the \emph{balance} between the two GP signals
is what separates effective from ineffective acquisition strategies.
Both predictions are borne out by the Uncertainty–Random equivalence.

\subsection{EI and UCB Are Interchangeable at This Budget}
\label{sec:ei-ucb}

Within the balanced-strategy class, EI and UCB are statistically
indistinguishable (Nemenyi $p = 0.404$), despite different theoretical
properties: UCB has finite-time regret guarantees under appropriate
$\beta$ schedules, whereas EI has no such guarantee.
The per-problem medians echo this: UCB leads on four of six problems, EI
on two, with no consistent pattern by problem type or dimension.

De~Ath et al.~\citeyear{deAth2019greed} provide a mechanistic
explanation: both EI and UCB lie on the exploration–exploitation Pareto
front, so they explore at similar rates; any residual difference is
smaller than the problem-to-problem and seed-to-seed variance at a
20-iteration budget.
Shahriari et al.~\citeyear{shahriari2016taking} reinforce this with a
broader observation: the surrogate model's adequacy typically dominates
the fine-grained acquisition choice.
The EI–UCB equivalence is therefore not a sample-size failure but the
expected outcome of using two well-designed balanced strategies on the
same adequate surrogate.

\subsection{GP Normalization as a First-Class Methodological Result}
\label{sec:normalization}

The GP normalization result warrants discussion beyond a footnote in the
Methods section.
The initial experiment on Piston (7D) showed EI and UCB performing at
random-search level—plausible only if the surrogate were failing.
Piston's inputs span approximately seven orders of magnitude—initial gas
volume $V_0 \in [0.002,\,0.010]\,\mathrm{m}^3$ alongside ambient
pressure $P_0 \in [90\,000,\,110\,000]\,\mathrm{Pa}$—so an
un-normalized GP cannot learn meaningful per-dimension lengthscales.
Adding \texttt{Normalize} and \texttt{Standardize} recovered ${\approx}20\times$
improvement for EI and ${\approx}200\times$ for UCB, while Random—
which bypasses the GP—remained unchanged, isolating the fault to the
surrogate.

This finding directly instantiates Shahriari et al.'s
\citeyear{shahriari2016taking} claim that surrogate adequacy dominates
acquisition choice: the problem was not that EI or UCB was the wrong
acquisition function, but that the surrogate itself was miscalibrated.
It also underscores a practical lesson for engineers deploying BO on
real systems: physical inputs with heterogeneous scales require explicit
normalization to the known search bounds, not data-adaptive bounds that
drift as observations accumulate.

\subsection{Limitations}
\label{sec:limitations}

\paragraph{Budget.}
The 20-iteration budget is appropriate for expensive black-box
functions~\citep{frazier2018tutorial}, but it is tight relative to the
10D Ackley search space.
At this budget, no strategy explores more than a small fraction of
$[-5, 5]^{10}$, which suppresses the EI–Random advantage from $47\times$
on 2D Branin to $2\times$ on 10D Ackley and limits conclusions about
asymptotic performance.
Longer-horizon experiments would be needed to determine whether
balanced strategies accumulate a growing advantage in high dimensions
or plateau.

\paragraph{Fixed $\beta$ for UCB.}
I use a fixed $\beta = 2$ throughout, which does not follow the
theory-motivated schedule $\beta_t \propto \log t$ under which GP-UCB's
regret guarantees hold~\citep{frazier2018tutorial}.
A theory-consistent $\beta$ schedule may widen or narrow the EI–UCB gap;
I cannot rule out this effect from the current experiment.

\paragraph{Sequential-only evaluation.}
All results use sequential ($q = 1$) BO.
Batch acquisition ($q > 1$, e.g.\ qEI) is important in practice when
multiple evaluations can run in parallel, and the relative performance of
strategies may differ in the batch setting.

\paragraph{Synthetic and physics-emulator problems only.}
The benchmark includes six analytic or closed-form problems.
Real engineering simulations—such as CFD or aerodynamic design codes—may
exhibit noise, discontinuities, or constraint violations not present
here.
Whether the two-clique finding persists on black-box simulation outputs
remains an open question.

% ============================================================
% §6  Conclusion
% ============================================================
\section{Conclusion}
\label{sec:conclusion}

I benchmarked four acquisition strategies—EI, UCB, Uncertainty, and
Random—on six problems spanning dimensions 2 to 10 and two problem
families (synthetic and engineering), using a Friedman–Nemenyi
non-parametric framework with $N = 60$ blocks.
The results converge on two findings.

First, the 240 BO runs partition cleanly into two equivalence classes:
balanced strategies \{EI, UCB\} are statistically superior to
pure-exploration strategies \{Uncertainty, Random\}, with all
cross-class Nemenyi $p$-values below $10^{-6}$ and an inter-class rank
gap of $1.62 \gg \mathrm{CD} = 0.61$.
The Uncertainty–Random tie ($p = 0.974$) is the sharpest statement of
this result: using a GP for exploration only, without any exploitation
of the posterior mean, provides no measurable advantage over blind
random search.

Second, GP normalization to the known problem bounds is a
first-class engineering requirement on multi-scale problems.
Without it, BO on Piston degraded to random-search performance;
restoring it recovered one to two orders of magnitude in final regret.
This confirms the theoretical expectation that surrogate adequacy
dominates the acquisition-function choice and highlights a pitfall
specific to physical systems with heterogeneous input scales.
